% !TEX root = main.tex
\chapter{Tìm hiểu và nghiên cứu thị trường}\label{chap:thị-truong}

\section{Thực trạng và vấn đề cốt lõi trong hoạt động marketing của
SME ngành F\&B tại Việt Nam}\label{sec:thuc-trang}

Thị trường F\&B tại Việt Nam có tính cạnh tranh rất cao, với sự bùng nổ
của các nền tảng mạng xã hội, trong đó \textbf{Facebook} đóng vai trò
then chốt trong việc thu hút và giữ chân khách hàng. Tuy nhiên, các chủ
doanh nghiệp F\&B vừa và nhỏ (SME) đang đối mặt với nhiều thách thức
trong hoạt động marketing:

\begin{itemize}
\item
  \textbf{Thiếu kiến thức và kỹ năng chuyên môn:} Hầu hết chủ quán là
  chuyên gia về ẩm thực và dịch vụ, nhưng lại thiếu kiến thức nền tảng
  về marketing số, sáng tạo nội dung và tối ưu hóa quảng cáo.
\item
  \textbf{Hạn chế về thời gian và nguồn lực:} Lịch trình vận hành quán
  ăn bận rộn khiến họ không có đủ thời gian để duy trì sự nhất quán
  trong việc tạo và đăng tải nội dung hàng ngày.
\item
  \textbf{Nội dung thiếu chuyên nghiệp và ngẫu hứng:} Các phương pháp
  hiện tại chủ yếu dựa vào chụp ảnh bằng điện thoại không qua chỉnh sửa
  chuyên nghiệp hoặc đăng bài theo cảm hứng tức thời, dẫn đến chất lượng
  nội dung thấp, không có chiến lược rõ ràng.
\item
  \textbf{Chi phí Agency đắt đỏ:} Việc thuê các agency marketing chuyên
  nghiệp thường vượt quá ngân sách của các SME, khiến họ phải tự xoay sở
  với các công cụ rời rạc.
\end{itemize}

Những vấn đề này dẫn đến việc các SME F\&B bỏ lỡ cơ hội tiếp cận khách
hàng tiềm năng, giảm hiệu quả kinh doanh và gặp khó khăn trong việc xây
dựng thương hiệu bền vững. Thực trạng này được phản ánh rõ trong các báo
cáo ngành gần đây. Theo
e-Conomy SEA 2025 của Google, Temasek và Bain \& Company
\tlcite{ref:economy2025}, riêng phân khúc vận tải và giao đồ ăn tại Việt Nam
đã đạt 5 tỷ USD năm 2025 với tốc độ tăng trưởng 20\% so với cùng kỳ ---
cho thấy quy mô thị trường F\&B số đang mở rộng nhanh chóng. Song song
đó, 81\% người dùng Việt Nam tương tác với AI hằng ngày, dẫn đầu khu
vực Đông Nam Á, với doanh thu ứng dụng tích hợp tính năng AI tăng 78\%
trong nửa đầu năm 2025 so với cùng kỳ năm trước \tlcite{ref:economy2025}. Tuy
nhiên, theo McKinsey Global Survey 2025 \tlcite{ref:mckinsey2025}, các doanh
nghiệp nhỏ có doanh thu dưới 100 triệu USD chỉ có 29\% đạt giai đoạn
triển khai AI ở quy mô, so với 49\% ở doanh nghiệp lớn --- khoảng cách
này phản ánh đúng thực trạng các SME ngành F\&B đang gặp phải: thị
trường tăng trưởng nhanh nhưng năng lực số hóa còn hạn chế, và giải pháp
AI tự động hóa marketing chuyên biệt là nhu cầu cấp thiết.

\section{Phân tích các sản phẩm đang có trên thị
trường}\label{sec:phan-tich-san-pham}

\subsection{Predis AI}\label{subsec:predis}

\begin{itemize}
\item
  Đặc điểm: Đây là công cụ AI Marketing AIO (All-in-One), tạo ra các
  bài đăng (cả hình ảnh/video), captions, hashtags từ một prompt duy
  nhất.
\item
  Phân tích:

  \begin{itemize}
  \item
    Điểm mạnh: Tự động hóa cao, tạo ra visual content nhanh, có tính
    năng resize cho các nền tảng (ví dụ định dạng video ngắn/Reels trên
    Facebook).
  \item
    Điểm yếu: Nội dung visual thường là stock footage hoặc AI generate,
    thiếu tính chân thực của quán ăn cụ thể; không tập trung đủ sâu vào
    thị trường F\&B Việt Nam; không có khả năng tùy chỉnh theo dữ liệu
    thương hiệu nội bộ của từng doanh nghiệp.
  \end{itemize}
\end{itemize}

Một điểm hạn chế căn bản của Predis AI là thiếu khả năng tích hợp dữ
liệu thương hiệu nội bộ: người dùng không thể cung cấp cho hệ thống
thông tin về menu thực tế, giá cả địa phương, hay cách pha chế đặc
trưng của quán. Hệ quả là nội dung sinh ra thường mang phong cách chung
chung, thiếu chi tiết đặc thù phân biệt thương hiệu với đối thủ cạnh
tranh --- đây là lý do tại sao chủ quán F\&B khó dùng Predis AI để xây
dựng bản sắc thương hiệu lâu dài mà không phải viết lại phần lớn nội
dung sau khi AI sinh.

\subsection{Buffer}\label{subsec:buffer}

\begin{itemize}
\item
  \textbf{Đặc điểm:} Buffer là nền tảng quản lý mạng xã hội cho phép
  người dùng lên lịch, đăng bài và theo dõi hiệu suất trên nhiều kênh
  (Facebook, Instagram, X, LinkedIn\ldots) từ một giao diện duy nhất.
  Buffer tập trung vào các doanh nghiệp nhỏ và nhà sáng tạo nội dung
  cần quản lý đa kênh một cách có tổ chức.
\item
  \textbf{Phân tích:}

  \begin{itemize}
  \item
    \textit{Điểm mạnh:} Giao diện đơn giản, dễ dùng ngay cả với người
    không chuyên kỹ thuật; hỗ trợ lên lịch đăng bài tự động và phân
    tích hiệu suất bài viết; gói miễn phí phù hợp với SME ngân sách
    thấp.
  \item
    \textit{Điểm yếu:} Buffer \textbf{không có khả năng sinh nội dung}
    --- người dùng phải tự viết toàn bộ caption và chuẩn bị hình ảnh
    trước khi đăng; không lưu trữ hay tận dụng dữ liệu thương hiệu nội
    bộ; không có cơ chế đánh giá chất lượng nội dung trước khi đăng
    tải. Buffer giải quyết bài toán \emph{phân phối} nội dung, nhưng
    bỏ ngỏ hoàn toàn bài toán \emph{sáng tạo} nội dung.
  \end{itemize}
\end{itemize}

\subsection{Invideo AI}\label{subsec:invideo}

\begin{itemize}
\item
  Đặc điểm: Chuyên biến chữ viết (kịch bản) thành video hoàn chỉnh bằng
  stock footage.
\item
  Phân tích:

  \begin{itemize}
  \item
    Điểm mạnh: Tạo video nhanh chóng, chất lượng khá tốt, dễ sử dụng.
  \item
    Điểm yếu: Phụ thuộc vào kịch bản đầu vào; video dùng stock footage
    nên không có hình ảnh món ăn thật của quán; đòi hỏi người dùng phải
    tự viết kịch bản trước.
  \end{itemize}
\end{itemize}

\subsection{Canva Magic Studio}\label{subsec:canva}

\begin{itemize}
\item
  Đặc điểm: Nền tảng thiết kế cực kỳ phổ biến với SME, giờ có thêm AI
  hỗ trợ viết lách và tạo ảnh.
\item
  Phân tích:

  \begin{itemize}
  \item
    Điểm mạnh: Dễ dùng, chủ quán F\&B đã quen sử dụng, tích hợp thiết
    kế và AI viết caption cơ bản.
  \item
    Điểm yếu: AI chỉ là tính năng phụ trợ (assistant), không tự động hóa
    toàn bộ chiến dịch; không có khả năng lưu trữ dữ liệu thương hiệu
    hay tự động đăng bài theo lịch.
  \end{itemize}
\end{itemize}

\subsection{AutoNuoi (autonuoi.com)}\label{subsec:autonuoi}

\begin{itemize}
\item
  \textbf{Đặc điểm:} AutoNuoi là phần mềm tự động hóa Facebook của Việt
  Nam, tập trung vào đăng bài hàng loạt lên nhiều nhóm, tự động tăng
  lượt thích/bình luận và quản lý đồng thời nhiều tài khoản trên cùng
  một nền tảng.
\item
  \textbf{Phân tích:}

  \begin{itemize}
  \item
    \textit{Điểm mạnh:} Công cụ nội địa quen thuộc với thị trường Việt
    Nam; tự động hóa các thao tác thủ công lặp lại trên Facebook; chi
    phí thấp so với agency.
  \item
    \textit{Điểm yếu:} Hoàn toàn không có khả năng sinh nội dung ---
    người dùng phải tự chuẩn bị bài viết trước khi đăng, bài toán thiếu
    ý tưởng sáng tạo không được giải quyết. Không có cơ chế kiểm soát
    chất lượng nội dung, không lưu trữ hay tận dụng dữ liệu thương
    hiệu. Quan trọng hơn, phương thức hoạt động (đăng hàng loạt, tăng
    tương tác ảo, nuôi nick) có nguy cơ vi phạm chính sách của Facebook,
    tạo rủi ro khóa tài khoản --- điều đặc biệt nguy hiểm với SME F\&B
    đang xây dựng thương hiệu lâu dài.
  \end{itemize}
\end{itemize}

\textit{Lưu ý về thời điểm:} AutoNuoi xuất hiện sau khi đề tài này được
khởi động; tuy nhiên sự tồn tại của nó xác nhận nhu cầu thị trường thực
sự tồn tại đối với công cụ marketing Facebook tự động tại Việt Nam, đồng
thời cho thấy khoảng trống mà các giải pháp hiện tại chưa lấp đầy: tự
động hóa có kiểm soát chất lượng và có khả năng sinh nội dung thông minh.

\subsection{Tổng hợp so sánh}

Bảng~\ref{tab:so-sanh-cong-cu} tổng hợp năm công cụ theo các tiêu chí
cốt lõi liên quan đến bài toán của đề tài. Ký hiệu: \checkmark\ = có
đầy đủ, $\sim$ = có một phần / hạn chế, $\times$ = không có.

\begin{table}[htbp]
\centering
\caption{So sánh các công cụ/sản phẩm hiện có trên thị trường.}
\label{tab:so-sanh-cong-cu}
\small
\resizebox{\textwidth}{!}{%
\begin{tabular}{lccccc}
\toprule
\textbf{Tiêu chí} & \textbf{Predis AI} & \textbf{Buffer} & \textbf{Invideo AI} & \textbf{Canva} & \textbf{AutoNuoi} \\
\midrule
Sinh nội dung văn bản tự động    & \checkmark & $\times$   & $\sim$     & $\sim$     & $\times$ \\
Lưu trữ dữ liệu thương hiệu      & $\times$   & $\times$   & $\times$   & $\times$   & $\times$ \\
Lên lịch \& đăng bài tự động     & \checkmark & \checkmark & $\times$   & $\times$   & \checkmark \\
Đánh giá chất lượng nội dung     & $\times$   & $\times$   & $\times$   & $\times$   & $\times$ \\
Phù hợp thị trường Việt Nam      & $\sim$     & $\sim$     & $\sim$     & $\sim$     & \checkmark \\
Bước duyệt nội dung (Approval)   & $\times$   & $\times$   & $\times$   & $\times$   & $\times$ \\
Rủi ro vi phạm chính sách nền tảng & $\times$ & $\times$   & $\times$   & $\times$   & \checkmark \\
\bottomrule
\end{tabular}}
\end{table}

Qua bảng so sánh, có thể thấy \textbf{không có công cụ nào} đồng thời
đáp ứng cả ba yêu cầu cốt lõi: (1) sinh nội dung marketing thông minh
dựa trên dữ liệu thương hiệu, (2) đánh giá chất lượng tự động trước khi
đăng, và (3) hỗ trợ bước duyệt nội dung của người dùng. Đây chính là
khoảng trống mà hệ thống AI Agent được đề xuất trong đề án này hướng
tới lấp đầy.

Điểm mấu chốt từ bảng so sánh là \emph{không có bất kỳ công cụ nào}
tích hợp đồng thời khả năng lưu trữ dữ liệu thương hiệu và cơ chế đánh
giá chất lượng trước khi đăng. Đây không phải ngẫu nhiên: hai tính năng
này đòi hỏi hiểu biết sâu về domain F\&B và khả năng kết hợp nhiều
nguồn thông tin (menu, brand guidelines, lịch sử bài đăng) thành một
bài viết nhất quán --- năng lực mà các công cụ hiện tại, dù mạnh về
thiết kế hay lên lịch, chưa hướng tới. Thêm vào đó, hàng rào ``rủi ro
vi phạm chính sách nền tảng'' mà AutoNuoi gây ra càng củng cố yêu cầu
cần một giải pháp vừa tự động hóa thông minh vừa tuân thủ nguyên tắc
nền tảng và bảo vệ thương hiệu dài hạn.

\section{Khảo sát nhu cầu chủ doanh
nghiệp}\label{sec:khao-sat}

\subsection{Lý do thực hiện}

Khảo sát thực tế được tiến hành nhằm:

\begin{itemize}
\item
  Hiểu thực trạng hoạt động marketing của các chủ doanh nghiệp F\&B tại
  Việt Nam, bao gồm cách họ sản xuất nội dung và quản lý các kênh mạng
  xã hội, đặc biệt là Facebook.
\item
  Nhận diện các vấn đề và khó khăn mà chủ quán F\&B gặp phải trong việc
  duy trì sự nhất quán, chất lượng nội dung và sự thiếu hụt kiến thức
  marketing chuyên môn.
\item
  Thu thập ý kiến đóng góp từ thực tiễn về mức độ sẵn sàng sử dụng AI
  để tự động hóa quy trình tạo và đăng nội dung marketing.
\item
  Nắm bắt các thói quen sử dụng mạng xã hội, công cụ hiện có và mong
  muốn cải thiện hiệu quả marketing, từ đó làm cơ sở thiết kế hệ thống
  AI Agent thân thiện với người dùng không chuyên IT.
\end{itemize}

\subsection{Kế hoạch và nội dung khảo sát}

Kế hoạch khảo sát là một bước quan trọng để đảm bảo thu thập dữ liệu
chính xác, phản ánh đúng thực trạng quản lý tri thức và thông tin tại
SME. Nội dung khảo sát được thiết kế nhằm đánh giá cách doanh nghiệp lưu
trữ, tổ chức và truy xuất thông tin, đồng thời xác định những nhu cầu và
mong muốn về giải pháp AI Agent.

\textbf{Đối tượng khảo sát:} Chủ sở hữu / Chủ quán (Owner); Quản lý nhà
hàng / Quản lý chi nhánh (Manager).

\textbf{Quy mô mẫu khảo sát:} 5 quán cà phê / đồ uống quy mô nhỏ và
1 nhà hàng ăn uống tại Việt Nam. Đây là khảo sát thăm dò định tính
(exploratory survey) nhằm định hướng thiết kế tính năng, không hướng
đến mức độ đại diện thống kê.

Nội dung phiếu khảo sát gồm các phần: thông tin cá nhân và doanh nghiệp;
thực trạng hoạt động marketing trên mạng xã hội; nhu cầu về giải pháp
AI Marketing; đánh giá mức độ cần thiết của các tính năng đề xuất (thang
1--5); đề xuất thêm.

\subsection{Kết quả khảo sát}

\textbf{Số lượng mẫu khảo sát:} 6 doanh nghiệp SME hoạt động trong ngành
Thực phẩm \& Đồ uống (F\&B) tại thị trường Việt Nam. Đây là khảo sát
thăm dò định tính (exploratory survey) nhằm định hướng thiết kế tính năng,
không mang mục tiêu thống kê đại diện.

\textbf{Phân loại đối tượng khảo sát:}

\begin{itemize}
\item
  \textbf{Phân theo ngành hàng F\&B:}

  \begin{itemize}
  \item
    83\% (5 mẫu) thuộc lĩnh vực Quán cà phê / Đồ uống quy mô nhỏ.
  \item
    17\% (1 mẫu) thuộc lĩnh vực Nhà hàng ăn uống.
  \end{itemize}
\item
  \textbf{Phân theo vị trí công việc:}

  \begin{itemize}
  \item
    83\% (5 mẫu) là Chủ sở hữu / Chủ quán.
  \item
    17\% (1 mẫu) là Quản lý nhà hàng.
  \end{itemize}
\end{itemize}

\textbf{Tóm tắt các vấn đề cốt lõi thường gặp:}
Kết quả khảo sát xác nhận rằng các chủ quán F\&B quy mô nhỏ có nhận thức
rõ tầm quan trọng của marketing trên mạng xã hội (\textbf{100\% sử dụng
Facebook}; một số đồng thời dùng thêm kênh khác) nhưng gặp rào cản lớn
trong việc triển khai hiệu quả.

\begin{itemize}
\item
  \textbf{Vấn đề về Thời gian \& Tính nhất quán:} Vấn đề \textbf{``Không
  có thời gian duy trì sự nhất quán (đăng bài đều đặn)''} được đánh giá
  là rào cản lớn nhất. Đa số người khảo sát dành dưới 3 giờ/tuần cho
  hoạt động này, cho thấy nhu cầu cấp thiết về tự động hóa.
\item
  \textbf{Vấn đề về Chất lượng và Sáng tạo Nội dung:} \textbf{``Thiếu ý
  tưởng nội dung sáng tạo''} và \textbf{``Chất lượng hình ảnh/video thấp''}
  là những khó khăn phổ biến, dẫn đến việc nội dung thường được làm
  ``ngẫu hứng'' hoặc dựa trên cảm hứng tìm kiếm thủ công, thiếu tính
  chuyên nghiệp.
\item
  \textbf{Giải pháp hiện tại:} Các công cụ hiện tại được sử dụng rời rạc
  (Excel, ghi chú thủ công) hoặc các phần mềm quản lý cơ bản, không hỗ
  trợ tính năng tạo nội dung hay tự động hóa đăng bài marketing.
\end{itemize}

\textbf{Đánh giá mức độ cần thiết các tính năng của AI Agent (Thang điểm
1--5):}

Bảng~\ref{tab:ket-qua-khao-sat} trình bày mức độ cần thiết của các tính
năng được đề xuất, dựa trên phản hồi của 6 đối tượng khảo sát.

\begin{table}[htbp]
\centering
\caption{Kết quả khảo sát: điểm trung bình các tính năng.}
\label{tab:ket-qua-khao-sat}
\small
\begin{tabular}{p{9.2cm} c}
\toprule
\textbf{Tính năng} & \textbf{Điểm trung bình} \\
\midrule
E. Tự động lên lịch và đăng bài lên Facebook Page & \textbf{4.8}/5 \\
C. AI tạo ra nội dung/caption với giọng văn chuyên nghiệp & 4.5/5 \\
B. AI phân tích cảm hứng thành ý tưởng marketing & 4.4/5 \\
G. Đề xuất các xu hướng (hashtag, âm nhạc hot) tại VN & 4.3/5 \\
A. Thu thập cảm hứng (video, blog) từ trình duyệt dễ dàng & 4.1/5 \\
D. AI chỉnh sửa/nâng cao chất lượng ảnh/video đầu vào & 3.9/5 \\
F. Chatbot trả lời câu hỏi về các nội dung đã lưu trữ & 3.5/5 \\
\bottomrule
\end{tabular}
\end{table}

\textbf{Đề xuất từ người được khảo sát:}

\begin{itemize}
\item
  Người dùng đề xuất giải pháp nên ưu tiên phát triển trên nền tảng
  \textbf{Mobile App} (do chủ quán thường xuyên di chuyển) hoặc tích hợp
  với các ứng dụng nhắn tin (Zalo, Messenger) để tiện sử dụng khi đang
  làm việc tại quán.
\item
  Yêu cầu bắt buộc phải có bước \textbf{duyệt (Approval step)} trước khi
  hệ thống AI tự động đăng bài lên các kênh mạng xã hội, nhằm kiểm soát
  rủi ro về thương hiệu.
\item
  Mong muốn AI có thể gợi ý \textbf{thời điểm vàng (prime time)} để đăng
  bài, tối ưu hóa lượng tương tác tại thị trường địa phương.
\item
  Mong muốn hệ thống \textbf{ghi nhớ phong cách thương hiệu} sau mỗi
  lần duyệt bài --- tức là càng dùng, hệ thống càng hiểu ``gu'' của
  quán và sinh nội dung sát hơn mà không cần chủ quán phải nhập lại
  yêu cầu từ đầu mỗi chiến dịch.
\item
  Một số chủ quán bày tỏ mong muốn hệ thống có khả năng \textbf{phân
  tích hiệu quả bài đăng} (lượt tiếp cận, tương tác) và từ đó tự
  điều chỉnh phong cách sinh nội dung trong các chiến dịch tiếp theo,
  hướng đến một vòng cải thiện không cần sự can thiệp thủ công liên tục.
\end{itemize}

\subsection{Kết luận khảo sát}

Từ kết quả khảo sát thực tế tại 6 doanh nghiệp SME ngành F\&B (bao gồm
các quán cà phê/đồ uống quy mô nhỏ và nhà hàng), có thể rút ra các kết
luận chính sau:

Vấn đề cốt lõi không phải là thiếu nhận thức, mà là thiếu khả năng thực
thi: các chủ quán đều hiểu tầm quan trọng của marketing trên Facebook;
rào cản lớn nhất là sự hạn chế về thời gian và nguồn lực để duy trì sự
nhất quán trong đăng bài hàng ngày.

Phương pháp thủ công không hiệu quả: việc lưu trữ và quản lý nội dung
qua các công cụ rời rạc (Excel, ghi chú điện thoại) đang gây lãng phí
thời gian và dẫn đến chất lượng nội dung không chuyên nghiệp, thiếu chiến
lược rõ ràng.

Nhu cầu về tự động hóa là cấp thiết nhất: tính năng tự động lên lịch và
đăng bài lên Facebook Page được đánh giá là tính năng quan trọng nhất
(điểm trung bình \textbf{4.8}/5), khẳng định nhu cầu tự động hóa cao độ
để giải phóng thời gian cho chủ quán F\&B.

Cần sự cân bằng giữa AI và kiểm soát con người: mặc dù mong muốn tự động
hóa cao, người dùng vẫn nhấn mạnh yêu cầu phải có bước duyệt nội dung
cuối cùng (Approval step) để đảm bảo an toàn thương hiệu và chất lượng
nội dung trước khi đăng tải công khai.

Ưu tiên phát triển nội dung marketing hơn quản lý tri thức nội bộ: các
tính năng liên quan đến sáng tạo nội dung (tạo caption chuyên nghiệp,
phân tích cảm hứng, đề xuất trend địa phương) được ưu tiên cao hơn các
tính năng truy xuất dữ liệu nội bộ.

Ngoài ra, kết quả khảo sát còn cho thấy một xu hướng đáng chú ý: các
chủ quán có kinh nghiệm sử dụng mạng xã hội lâu hơn có xu hướng đánh
giá cao hơn các tính năng liên quan đến phân tích và tối ưu nội dung
(Tính năng B, G), trong khi các chủ quán mới hơn ưu tiên tự động hóa
đăng bài (Tính năng E) và tạo nội dung cơ bản (Tính năng C). Điều này
gợi ý rằng hệ thống AI Agent cần được thiết kế với cơ chế \emph{tiệm
tiến}: đơn giản với người mới bắt đầu, và cung cấp thêm công cụ phân
tích cho người dùng đã quen với luồng làm việc cốt lõi.

Kết quả khảo sát xác nhận bối cảnh và vấn đề đã nêu trong phần đặt vấn
đề: các SME ngành F\&B hiện nay gặp khó khăn trong việc thực hiện
marketing hiệu quả, và việc phát triển một AI Agent chuyên biệt là một
giải pháp phù hợp và cấp thiết.

Chương 2 đã làm rõ ba luận điểm nền tảng: (1) các chủ quán F\&B nhận
thức được tầm quan trọng của marketing Facebook nhưng bị giới hạn bởi
thời gian và năng lực chuyên môn; (2) các công cụ hiện có trên thị
trường chưa giải quyết đồng thời cả ba yêu cầu cốt lõi là sinh nội dung
thông minh, đánh giá chất lượng và kiểm soát của người dùng; (3) người
dùng chấp nhận AI nhưng yêu cầu bắt buộc phải có bước duyệt nội dung
trước khi đăng. Ba luận điểm này dẫn dắt trực tiếp đến thiết kế hệ
thống được trình bày từ Chương~\ref{chap:ly-thuyet} trở đi.

\section{Sự cần thiết của quy trình kiểm duyệt nội dung}\label{sec:approval}

Bên cạnh nhu cầu tự động hóa, kết quả khảo sát cho thấy người dùng không
muốn giao toàn bộ quyền kiểm soát cho AI. Bước duyệt nội dung (Approval
step) trước khi đăng là yêu cầu bắt buộc được đề cập bởi tất cả các đối
tượng khảo sát, xuất phát từ lo ngại về rủi ro thương hiệu khi nội dung
không phù hợp được đăng tải công khai.

Đề tài hướng tới xây dựng \textbf{pipeline đánh giá chất lượng tự động
đa lớp} nhằm giảm gánh nặng kiểm duyệt thủ công lặp lại. Pipeline này
thực hiện lọc và chấm điểm nội dung theo quy tắc cứng, chỉ số xác suất
và đánh giá định tính trước khi trình lên người dùng phê duyệt. Cách tiếp
cận này kết hợp hiệu quả của tự động hóa với sự kiểm soát của con người,
đảm bảo chất lượng và tính nhất quán của nội dung được đăng tải.

Về mặt kỹ thuật, pipeline đánh giá đóng vai trò trung gian thông minh
giữa AI sinh nội dung và người dùng cuối: thay vì hiển thị thô mọi bài
AI tạo ra để người dùng tự phán xét, hệ thống chỉ trình lên những bài
vượt ngưỡng chất lượng tối thiểu. Điều này không chỉ tiết kiệm thời
gian duyệt bài mà còn giúp chủ quán tin tưởng hơn vào đầu ra của AI,
vì họ biết rằng những gì họ thấy đã qua một lớp kiểm tra khách quan.
Mối quan hệ giữa yêu cầu bước duyệt từ khảo sát và thiết kế pipeline
đánh giá trong hệ thống sẽ được trình bày chi tiết tại
Chương~\ref{chap:thuc-nghiem}.
